\section{Introduction}
\label{sec:intro}

A security operator watches a live camera feed and asks: \emph{``Did anyone enter the building in the last five minutes?''}
Someone wearing smart glasses asks: \emph{``What did I put in my bag?''}
Both questions target a video stream that is still running. Standard video question answering (VideoQA) cannot help here. It assumes the full video is available before inference starts.

Recent models like Video-ChatGPT~\cite{maaz2023videochatgpt}, VideoLLaVA~\cite{lin2023videollava}, and LLaVA-Next-Video~\cite{zhang2024llavanextvideo} score well on standard VideoQA benchmarks, but they share one assumption: the video is finite and fully available at query time.
Frame sampling, global temporal pooling, and bidirectional attention all depend on that assumption.
A live stream has no end. Feed one into these models and three problems surface:

\begin{enumerate}[leftmargin=*,itemsep=2pt]
    \item \textbf{Unbounded memory.} Encoding every frame from a continuous stream means memory and compute grow linearly. Hardware limits will stop the system.
    \item \textbf{Temporal ambiguity.} \emph{``What is happening?''} targets the present. \emph{``Was there a dog earlier?''} targets the past. Without explicit temporal reasoning, the model mixes these contexts and hallucinates.
    \item \textbf{Latency.} Re-processing the full stream for each query adds delay. Users expect sub-second responses.
\end{enumerate}

\textbf{\ours} solves all three. It is a streaming vision-language system for interactive Q\&A on live video, built from three components:

\begin{enumerate}[leftmargin=*,itemsep=2pt]
    \item \textbf{Semantic Keyframe Memory (SKM).}
    A fixed bank of $N$ entries stores the most informative frames, scored by visual novelty and temporal coverage. High-scoring frames enter; low-scoring stored frames get evicted.

    \item \textbf{Temporal Query Router (TQR).}
    The TQR classifies each question into one of three scopes---\emph{instantaneous}, \emph{recent}, or \emph{historical}---and hands only the matching memory slice to the answer stage.

    \item \textbf{Stream-Fused Generator (SFG).}
    BLIP~\cite{li2022blip} runs captioning and VQA on each filtered keyframe; Flan-T5~\cite{chung2022flanT5} then fuses those observations into a single response.
    Non-neural components run in under 2\,ms; on a GPU the full pipeline completes in under one second per query.
\end{enumerate}

For evaluation, we convert NExT-QA~\cite{xiao2021nextqa} and EgoSchema~\cite{mangalam2023egoschema} into streaming protocols: the evaluator issues each question at a specific timestamp, and the model sees no future frames. We also evaluate on OVO-Bench~\cite{li2025ovobench}, a recent benchmark explicitly designed for online video understanding, and on a streaming subset of Ego4D-NLQ~\cite{grauman2022ego4d}.
We introduce \emph{LiveQA-Bench}, 500 temporally-diverse questions spread across 50 simulated live streams of 5 to 30 minutes, covering past, present, and transitional events.

We compare against eight baselines, including recent streaming-specific models: Flash-VStream~\cite{zhang2024flashvstream}, VideoLLM-online~\cite{chen2024videollmonline}, and Dispider~\cite{he2024dispider}.
On LiveQA-Bench, \ours scores 39.7\% overall: 59.3\% on instant queries, 33.3\% on historical queries requiring long-range memory, and 16.7\% on recent queries.
Instant questions are easiest because the latest keyframe suffices; historical questions benefit from the SKM's importance-based retention.
Streaming evaluations on NExT-QA, EgoSchema, OVO-Bench, and Ego4D-NLQ, together with baseline comparisons, will be reported in the final version.

A live web demo, suitable for workshop presentation, lets users ask questions about their own webcam feed in real time.
